\documentclass[12pt,a4paper]{article}

\usepackage[T1]{fontenc}
\usepackage{lmodern}
\usepackage[margin=0pt]{geometry}
\usepackage{xcolor}
\usepackage{tikz}
\usepackage{microtype}
\usepackage{amsmath}
\usepackage{amssymb}
\usepackage{listings}
\usepackage{graphicx}
\usepackage{etoolbox}
\usepackage[hidelinks,linktoc=all]{hyperref}

% Give every manually added table-of-contents entry its own link target.
\pretocmd{\addcontentsline}{\phantomsection}{}{}

% Cover palette
\definecolor{accentgold}{HTML}{d4a94c}
\definecolor{coverdark}{HTML}{0d1117}
\definecolor{coverlight}{HTML}{21262d}
\definecolor{subtlewhite}{HTML}{c9d1d9}
\definecolor{dimwhite}{HTML}{8b949e}
\definecolor{sectionnavy}{HTML}{1e3a5f}

\lstset{
  language=C++,
  basicstyle=\ttfamily\footnotesize,
  keywordstyle=\bfseries\color{sectionnavy},
  commentstyle=\itshape\color{dimwhite},
  frame=single,
  rulecolor=\color{black!25},
  backgroundcolor=\color{black!2},
  breaklines=true,
  showstringspaces=false,
  columns=fullflexible,
  keepspaces=true
}

% Keep code blocks intact; if one does not fit, move it to the next page.
\BeforeBeginEnvironment{lstlisting}{%
  \par\noindent\begin{minipage}{\linewidth}%
}
\AfterEndEnvironment{lstlisting}{%
  \end{minipage}\par
}

\pagestyle{empty}

\begin{document}

% ----------------------- COVER PAGE -----------------------
\pagenumbering{gobble}
\pagecolor{coverdark}

\begin{tikzpicture}[remember picture, overlay]

  % ===== BACKGROUND LAYERS =====
  % Subtle top-to-bottom gradient
  \shade[top color=coverlight, bottom color=coverdark, opacity=0.6]
    (current page.south west) rectangle (current page.north east);



  % ===== MAIN TEXT CONTENT =====

  % — top accent line —
  \draw[accentgold, line width=0.8pt, opacity=0.4]
    ([yshift=-3.0cm, xshift=-4cm]current page.north)
    -- ++(8cm, 0);

  % — subtitle —
  \node[anchor=center, text=accentgold, font=\fontsize{11}{14}\selectfont]
    at ([yshift=-3.7cm]current page.north)
    {\textls[200]{THE COMPLETE SOLUTIONS MANUAL}};

  % — second accent line —
  \draw[accentgold, line width=0.4pt, opacity=0.25]
    ([yshift=-4.3cm, xshift=-3cm]current page.north)
    -- ++(6cm, 0);

  % — main title (two lines) —
  \node[anchor=center, text=white,
        font=\fontsize{36}{44}\selectfont\bfseries,
        align=center]
    at ([yshift=-7.3cm]current page.north)
    {Data Structures \&\\[6pt] Algorithm Analysis in C++};

  % — textbook author —
  \node[anchor=center, text=dimwhite,
        font=\fontsize{14}{18}\selectfont\itshape]
    at ([yshift=-10.0cm]current page.north)
    {by Mark Allen Weiss};

  % — decorative diamond —
  \node[anchor=center, text=accentgold, opacity=0.5,
        font=\fontsize{10}{10}\selectfont]
    at ([yshift=-11.3cm]current page.north)
    {$\diamond$};

  % — chapter block (no frame) —
  \node[anchor=center, text=accentgold,
        font=\fontsize{22}{26}\selectfont\bfseries]
    at ([yshift=-13.3cm]current page.north)
    {Chapter 2};
  \node[anchor=center, text=subtlewhite,
        font=\fontsize{18}{22}\selectfont]
    at ([yshift=-14.4cm]current page.north)
    {Algorithm Analysis};

  % — asymptotic-complexity motif —
  \begin{scope}[
      shift={(current page.south)},
      yshift=8.4cm,
      x=1.12cm,
      y=1.12cm,
      scale=1.5,
      transform shape
    ]
    % Large, quiet Big-O watermark
    \node[text=accentgold, opacity=0.045,
          font=\fontsize{92}{92}\selectfont\bfseries]
      at (0,0.35) {$\mathcal{O}$};

    % Axes
    \draw[->, subtlewhite, opacity=0.42, line width=0.65pt]
      (-3.5,-1.15) -- (3.65,-1.15)
      node[below left=1pt, opacity=0.85,
           font=\scriptsize\itshape] {$n$};
    \draw[->, subtlewhite, opacity=0.42, line width=0.65pt]
      (-3.5,-1.15) -- (-3.5,1.85)
      node[above right=1pt, opacity=0.85,
           font=\scriptsize\itshape] {$T(n)$};

    % Growth-rate curves
    \begin{scope}[shift={(-3.5,-1.15)}]
      \draw[subtlewhite, opacity=0.45, line width=1.15pt,
            domain=0.12:7.0, samples=70, smooth]
        plot (\x,{0.34*ln(1+1.25*\x)});
      \node[anchor=west, text=subtlewhite, opacity=0.62,
            font=\scriptsize]
        at (5.55,0.48) {$\log n$};

      \draw[accentgold, opacity=0.62, line width=1.25pt,
            domain=0:6.6, samples=2]
        plot (\x,{0.22*\x});
      \node[anchor=west, text=accentgold, opacity=0.78,
            font=\scriptsize]
        at (5.55,0.98) {$n$};

      \draw[accentgold, opacity=0.82, line width=1.45pt,
            domain=0.12:6.17, samples=70, smooth]
        plot (\x,{0.16*\x*ln(1+\x)});
      \node[anchor=west, text=accentgold, opacity=0.9,
            font=\scriptsize]
        at (4.62,1.82) {$n\log n$};

      \draw[accentgold, opacity=0.95, line width=1.55pt,
            domain=0:5.15, samples=70, smooth]
        plot (\x,{0.105*\x*\x});
      \node[anchor=west, text=accentgold,
            font=\scriptsize\bfseries]
        at (4.18,2.30) {$n^2$};

    \end{scope}
  \end{scope}

  % — bottom section —
  % Author name
  \node[anchor=center, text=subtlewhite,
        font=\fontsize{16}{20}\selectfont]
    at ([yshift=4.5cm]current page.south)
    {Ahmed Amine Azzouzi};

  % Bottom accent line
  \draw[accentgold, line width=0.6pt, opacity=0.3]
    ([yshift=3.5cm, xshift=-2.5cm]current page.south)
    -- ++(5cm, 0);

  % Small year tag
  \node[anchor=center, text=dimwhite, opacity=0.5,
        font=\fontsize{9}{12}\selectfont]
    at ([yshift=2.8cm]current page.south)
    {\textls[150]{2026}};

\end{tikzpicture}

% Supply a real page box for the overlay.
\null

% -------------------- TABLE OF CONTENTS --------------------
\clearpage
\newgeometry{margin=2.5cm}
\pagecolor{white}
\color{black}
\pagestyle{plain}
\pagenumbering{roman}
\setcounter{page}{1}

\tableofcontents

% ------------------------- EXERCISES ------------------------
\clearpage
\pagenumbering{arabic}
\setcounter{page}{1}

\section*{Exercises}
\addcontentsline{toc}{section}{Exercises}
\noindent
{\color{accentgold}\rule{\linewidth}{0.7pt}}
\vspace{1.2em}

\subsection*{\color{sectionnavy}Exercise 2.1}
\addcontentsline{toc}{subsection}{Exercise 2.1: Ordering asymptotic growth rates}

\noindent
Order the following functions by growth rate:
$N$, $\sqrt{N}$, $N^{1.5}$, $N^2$, $N\log N$,
$N\log\log N$, $N\log^2 N$, $N\log(N^2)$,
$2/N$, $2^N$, $2^{N/2}$, $37$, $N^2\log N$, and $N^3$.
Indicate which functions grow at the same rate.

\subsubsection*{\color{sectionnavy}Solution}

As $N\to\infty$, the functions are ordered from slowest to fastest as
follows, where $f\prec g$ means that $f=o(g)$, and $f\asymp g$ means that $f=\Theta(g)$:
\[
\begin{aligned}
\frac{2}{N}
&\prec 37
 \prec \sqrt{N}
 \prec N
 \prec N\log\log N \\
&\prec N\log N
 \asymp N\log(N^2)
 \prec N(\log N)^2
 \prec N^{1.5} \\
&\prec N^2
 \prec N^2\log N
 \prec N^3
 \prec 2^{N/2}
 \prec 2^N.
\end{aligned}
\]

The only functions in the list that grow at the same asymptotic rate are
$N\log N$ and $N\log(N^2)$, since
\[
N\log(N^2)=2N\log N=\Theta(N\log N).
\]

\subsection*{\color{sectionnavy}Exercise 2.2}
\addcontentsline{toc}{subsection}{Exercise 2.2: Properties of Big-Oh bounds}

\noindent
Suppose $T_1(N)=O(f(N))$ and $T_2(N)=O(f(N))$. Which of the following
are true?

\begin{enumerate}
  \renewcommand{\labelenumi}{\alph{enumi}.}
  \item $T_1(N)+T_2(N)=O(f(N))$
  \item $T_1(N)-T_2(N)=o(f(N))$
  \item $\dfrac{T_1(N)}{T_2(N)}=O(1)$
  \item $T_1(N)=O(T_2(N))$
\end{enumerate}

\subsubsection*{\color{sectionnavy}Solution}

\begin{enumerate}
  \renewcommand{\labelenumi}{\alph{enumi}.}

  \item \textbf{True.}
  If $T_1(N)\leq c_1f(N)$ and $T_2(N)\leq c_2f(N)$ for sufficiently
  large $N$, then
  \[
    T_1(N)+T_2(N)\leq (c_1+c_2)f(N).
  \]

  \item \textbf{False.}
  For example, let $T_1(N)=2f(N)$ and $T_2(N)=f(N)$. Their difference
  is $f(N)$, which is not $o(f(N))$.

  \item \textbf{False.}
  Both functions can be bounded above by $f(N)$ while having very
  different sizes. For example, if $f(N)=N$, $T_1(N)=N$, and
  $T_2(N)=1$, then $T_1(N)/T_2(N)=N$, which is not $O(1)$.

  \item \textbf{False.}
  Using the same example, $T_1(N)=N$ is not $O(T_2(N))=O(1)$.
\end{enumerate}

Therefore, only statement~\textbf{(a)} is necessarily true.

\subsection*{\color{sectionnavy}Exercise 2.3}
\addcontentsline{toc}{subsection}{Exercise 2.3: Comparing two growth rates}

\noindent
Which function grows faster: $N\log N$ or
$N^{1+\varepsilon/\sqrt{\log N}}$, $\varepsilon>0$?

\subsubsection*{\color{sectionnavy}Solution}

Consider their ratio:
\[
\frac{N^{1+\varepsilon/\sqrt{\log N}}}{N\log N}
=
\frac{N^{\varepsilon/\sqrt{\log N}}}{\log N}.
\]
Using natural logarithms (a different fixed base changes only a positive
constant), this becomes
\[
\frac{e^{\varepsilon\sqrt{\log N}}}{\log N}.
\]
Its logarithm is
\[
\varepsilon\sqrt{\log N}-\log\log N,
\]
which tends to infinity because $\sqrt{\log N}$ grows faster than
$\log\log N$. Therefore,
\[
N^{1+\varepsilon/\sqrt{\log N}}
=\omega(N\log N),
\]
so $N^{1+\varepsilon/\sqrt{\log N}}$ grows faster.

\subsection*{\color{sectionnavy}Exercise 2.4}
\addcontentsline{toc}{subsection}{Exercise 2.4: Logarithms versus linear growth}

\noindent
Prove that for any constant $k$, $\log^k N=o(N)$.

\subsubsection*{\color{sectionnavy}Solution}

Let $x=\log N$. The base of the logarithm affects only a constant
factor, so we may use the natural logarithm. Then $N=e^x$, and
\[
\frac{\log^k N}{N}=\frac{x^k}{e^x}.
\]
Choose an integer $m>k$. From the power-series expansion of $e^x$,
for $x>0$ we have
\[
e^x\geq \frac{x^m}{m!}.
\]
Consequently,
\[
0\leq \frac{x^k}{e^x}
\leq \frac{m!}{x^{m-k}}
\longrightarrow 0
\qquad\text{as }x\to\infty.
\]
Therefore,
\[
\lim_{N\to\infty}\frac{\log^k N}{N}=0,
\]
which proves that $\log^k N=o(N)$ for every fixed constant $k$.

\subsection*{\color{sectionnavy}Exercise 2.5}
\addcontentsline{toc}{subsection}{Exercise 2.5: Incomparable asymptotic functions}

\noindent
Find two functions $f(N)$ and $g(N)$ such that neither
$f(N)=O(g(N))$ nor $g(N)=O(f(N))$.

\subsubsection*{\color{sectionnavy}Solution}

Define
\[
f(N)=
\begin{cases}
N,   & \text{if $N$ is even},\\
N^2, & \text{if $N$ is odd},
\end{cases}
\qquad
g(N)=
\begin{cases}
N^2, & \text{if $N$ is even},\\
N,   & \text{if $N$ is odd}.
\end{cases}
\]
For odd $N$,
\[
\frac{f(N)}{g(N)}=N,
\]
which is unbounded. Hence $f(N)\neq O(g(N))$. On the other hand, for
even $N$,
\[
\frac{g(N)}{f(N)}=N,
\]
which is also unbounded. Hence $g(N)\neq O(f(N))$. Therefore, neither
function asymptotically dominates the other.

\subsection*{\color{sectionnavy}Exercise 2.6}
\addcontentsline{toc}{subsection}{Exercise 2.6: A repeatedly squared fine}

\noindent
In a recent court case, a judge cited a city for contempt and ordered a
fine of \$2 for the first day. Each subsequent day, until the city
followed the judge's order, the fine was squared (i.e., the fine
progressed as follows: \$2, \$4, \$16, \$256, \$65{,}536, \ldots).

\begin{enumerate}
  \renewcommand{\labelenumi}{\alph{enumi}.}
  \item What would be the fine on day $N$?
  \item How many days would it take for the fine to reach $D$ dollars
  (a Big-Oh answer will do)?
\end{enumerate}

\subsubsection*{\color{sectionnavy}Solution}

Let $F_N$ be the fine on day $N$. The rule gives
\[
F_1=2,
\qquad
F_N=F_{N-1}^2.
\]

\begin{enumerate}
  \renewcommand{\labelenumi}{\alph{enumi}.}

  \item The exponent of $2$ doubles each day:
  \[
  F_1=2^1,\quad F_2=2^2,\quad F_3=2^4,\quad F_4=2^8.
  \]
  Therefore, the fine on day $N$ is
  \[
  \boxed{F_N=2^{2^{N-1}}}.
  \]

  \item To find the first day on which the fine is at least $D$, solve
  \[
  2^{2^{N-1}}\geq D.
  \]
  Taking the base-$2$ logarithm twice gives
  \[
  N-1\geq \log_2\log_2 D.
  \]
  Thus, for $D\geq2$, the first such day is
  \[
  \boxed{N=1+\left\lceil\log_2\log_2 D\right\rceil},
  \]
  and the number of days is $\Theta(\log\log D)$.
\end{enumerate}

\clearpage
\subsection*{\color{sectionnavy}Exercise 2.7}
\addcontentsline{toc}{subsection}{Exercise 2.7: Analyzing nested-loop fragments}

\noindent
For each of the following six program fragments:

\begin{enumerate}
  \renewcommand{\labelenumi}{\alph{enumi}.}
  \item Give an analysis of the running time (Big-Oh will do).
  \item Implement the code in the language of your choice, and give the
  running time for several values of $N$.
  \item Compare your analysis with the actual running times.
\end{enumerate}

\begin{lstlisting}
// (1)
sum = 0;
for (i = 0; i < n; ++i)
    ++sum;

// (2)
sum = 0;
for (i = 0; i < n; ++i)
    for (j = 0; j < n; ++j)
        ++sum;

// (3)
sum = 0;
for (i = 0; i < n; ++i)
    for (j = 0; j < n * n; ++j)
        ++sum;

// (4)
sum = 0;
for (i = 0; i < n; ++i)
    for (j = 0; j < i; ++j)
        ++sum;

// (5)
sum = 0;
for (i = 0; i < n; ++i)
    for (j = 0; j < i * i; ++j)
        for (k = 0; k < j; ++k)
            ++sum;

// (6)
sum = 0;
for (i = 1; i < n; ++i)
    for (j = 1; j < i * i; ++j)
        if (j % i == 0)
            for (k = 0; k < j; ++k)
                ++sum;
\end{lstlisting}

\subsubsection*{\color{sectionnavy}Solution}

\paragraph{Running-time analysis.}
Counting executions of the innermost statement gives:

\begin{enumerate}
  \item[(1)] The loop executes $N$ times, so the running time is
  $\Theta(N)$.

  \item[(2)] Each of the two loops executes $N$ times, giving
  $N^2$ increments and a running time of $\Theta(N^2)$.

  \item[(3)] The inner loop executes $N^2$ times for each of $N$ outer
  iterations. Thus the running time is $\Theta(N^3)$.

  \item[(4)] The number of increments is
  \[
    \sum_{i=0}^{N-1}i=\frac{N(N-1)}{2},
  \]
  so the running time is $\Theta(N^2)$.

  \item[(5)] For a fixed $i$, the innermost statement executes
  \[
    \sum_{j=0}^{i^2-1}j=\frac{i^2(i^2-1)}{2}=\Theta(i^4)
  \]
  times. Summing over $i$ gives
  $\Theta(\sum_{i=0}^{N-1}i^4)=\Theta(N^5)$.

  \item[(6)] For a fixed $i$, the condition succeeds when $j=mi$ for
  $m=1,\ldots,i-1$. The innermost statement therefore executes
  \[
    \sum_{m=1}^{i-1}mi
    =\frac{i^2(i-1)}{2}
    =\Theta(i^3)
  \]
  times. Summing over $i$ gives $\Theta(N^4)$. The $\Theta(N^3)$
  condition checks do not change this bound.
\end{enumerate}

Hence the six bounds are
\[
\boxed{\Theta(N),\ \Theta(N^2),\ \Theta(N^3),\
       \Theta(N^2),\ \Theta(N^5),\ \Theta(N^4)}.
\]

\paragraph{C++ implementation.}
The following functions were timed separately. A volatile accumulator is
used so that an optimizing compiler cannot remove the loops.

\begin{lstlisting}
#include <cstdint>

volatile std::uint64_t sum;

void fragment1(int n) {
    sum = 0;
    for (int i = 0; i < n; ++i) ++sum;
}

void fragment2(int n) {
    sum = 0;
    for (int i = 0; i < n; ++i)
        for (int j = 0; j < n; ++j) ++sum;
}

void fragment3(int n) {
    sum = 0;
    for (int i = 0; i < n; ++i)
        for (int j = 0; j < n * n; ++j) ++sum;
}

void fragment4(int n) {
    sum = 0;
    for (int i = 0; i < n; ++i)
        for (int j = 0; j < i; ++j) ++sum;
}

void fragment5(int n) {
    sum = 0;
    for (int i = 0; i < n; ++i)
        for (int j = 0; j < i * i; ++j)
            for (int k = 0; k < j; ++k) ++sum;
}

void fragment6(int n) {
    sum = 0;
    for (int i = 1; i < n; ++i)
        for (int j = 1; j < i * i; ++j)
            if (j % i == 0)
                for (int k = 0; k < j; ++k) ++sum;
}
\end{lstlisting}

\paragraph{Representative timings.}
The benchmark was compiled with Microsoft Visual C++ in Release x64 mode
and run on an Intel Core i7-13700HX. The following filtered measurements
are in microseconds. Absolute times depend on the machine; the scaling
ratios are the important part.

\begin{center}
\begin{tabular}{c|r|r|r|r|r|r}
$N$ & \textbf{F1} & \textbf{F2} & \textbf{F3} &
\textbf{F4} & \textbf{F5} & \textbf{F6} \\
\hline
10 & 0.0073 & 0.1450 & 1.0190 & 0.0438 & 14.4 & 2.1 \\
20 & 0.0117 & 0.1256 & 3.2540 & 0.3083 & 174.9 & 15.1 \\
30 & 0.0161 & 0.3765 & 7.1690 & 0.4341 & 640.9 & 79.7 \\
40 & 0.0190 & 0.8298 & 14.401 & 0.7059 & 3{,}782.3 & 113.6 \\
50 & 0.0261 & 0.6919 & 26.805 & 0.8927 & 8{,}979.9 & 224.6
\end{tabular}
\end{center}

\paragraph{Doubling-ratio comparison.}
For a running time of order $N^k$, doubling $N$ should multiply the time
by approximately $2^k$. Selected stable measurements from the scaling
benchmark are shown below.

\begin{center}
\resizebox{\textwidth}{!}{%
\begin{tabular}{c|c|c|c|c|c}
\textbf{Fragment} & \textbf{Complexity} & $N\to2N$ &
\textbf{Time ($\mu$s)} & \textbf{Observed} & \textbf{Expected} \\
\hline
1 & $\Theta(N)$   & 200k$\to$400k & 41.70$\to$84.00 & 2.01 & 2 \\
2 & $\Theta(N^2)$ & 4k$\to$8k     & 3{,}385.1$\to$13{,}556.1 & 4.00 & 4 \\
3 & $\Theta(N^3)$ & 400$\to$800   & 13{,}819.2$\to$108{,}560.0 & 7.86 & 8 \\
4 & $\Theta(N^2)$ & 4k$\to$8k     & 1{,}728.5$\to$6{,}988.2 & 4.04 & 4 \\
5 & $\Theta(N^5)$ & 40$\to$80     & 3{,}034.5$\to$84{,}129.0 & 27.72 & 32 \\
6 & $\Theta(N^4)$ & 40$\to$80     & 119.3$\to$1{,}936.3 & 16.23 & 16
\end{tabular}%
}
\end{center}

Fragments (1)--(4) and (6) closely match their theoretical doubling
factors. Fragment (5) is somewhat below the ideal factor of $32$ at the
tested sizes, but its much steeper scaling is consistent with
$\Theta(N^5)$. Timer resolution, caching, processor scheduling, and
compiler effects explain the remaining variation. Overall, the
i7-13700HX measurements support all six asymptotic bounds.

\clearpage
\subsection*{\color{sectionnavy}Exercise 2.8}
\addcontentsline{toc}{subsection}{Exercise 2.8: Random permutation algorithms}

\noindent
Suppose that we need to generate a random permutation of the first $N$
integers. For example, $\{4,3,1,5,2\}$ and $\{3,1,4,2,5\}$ are legal
permutations, whereas $\{5,4,1,2,1\}$ is not. Assume that
\texttt{randInt(i,j)} generates every integer between $i$ and $j$ with
equal probability. Consider the following algorithms.

\begin{enumerate}
  \item Fill $a[0]$ through $a[N-1]$. To fill $a[i]$, generate random
  integers until one is obtained that does not already occur in
  $a[0],\ldots,a[i-1]$.

  \item Proceed as in algorithm (1), but maintain a Boolean array
  \texttt{used}. When a value is inserted, mark it as used, allowing
  membership to be tested in constant time.

  \item First set $a[i]=i+1$, and then execute
\begin{lstlisting}
for (int i = 1; i < n; ++i)
    swap(a[i], a[randInt(0, i)]);
\end{lstlisting}
\end{enumerate}

\begin{enumerate}
  \renewcommand{\labelenumi}{\alph{enumi}.}
  \item Prove that all three algorithms produce only legal
  permutations and that all permutations are equally likely.
  \item Give as accurate a Big-Oh analysis as possible for the expected
  running time of each algorithm.
  \item Implement each algorithm, run it ten times for the requested
  values of $N$, and compute the average running times.
  \item Compare the analysis with the measured running times.
  \item Give the worst-case running time of each algorithm.
\end{enumerate}

\subsubsection*{\color{sectionnavy}Solution}

\paragraph{(a) Correctness and uniformity.}
Algorithms (1) and (2) accept only values not already present in the
array, so their outputs are legal permutations. When $i$ positions have
been filled, rejection sampling makes each of the $N-i$ unused values
equally likely to be selected next. Consequently, every particular
permutation has probability
\[
\frac1N\cdot\frac1{N-1}\cdots\frac11=\frac1{N!}.
\]

Algorithm (3) is the forward Fisher--Yates shuffle. After iteration $i$,
the prefix $a[0],\ldots,a[i]$ is a uniformly random permutation of its
$i+1$ values. This is true initially, and choosing the swap position
uniformly from $0$ through $i$ preserves the claim inductively. After
the final iteration, every one of the $N!$ permutations therefore has
probability $1/N!$.

\paragraph{(b) Expected running times.}
For algorithms (1) and (2), when $i$ positions have been filled, a
random draw succeeds with probability $(N-i)/N$. The expected number of
draws is therefore
\[
\frac{N}{N-i}.
\]

In algorithm (1), a successful membership search examines all $i$
existing entries. The expected number of failed draws is $i/(N-i)$,
and a failed search examines $(i+1)/2$ entries on average. Thus the
expected work at position $i$ is
\[
i+\frac{i}{N-i}\frac{i+1}{2}.
\]
Summing this expression gives
\[
\boxed{\mathbb{E}[T_1(N)]=\Theta(N^2\log N)}.
\]

Algorithm (2) tests \texttt{used} in constant time, so its expected
number of random draws is
\[
\sum_{i=0}^{N-1}\frac{N}{N-i}
=N\sum_{r=1}^{N}\frac1r
=NH_N.
\]
Therefore,
\[
\boxed{\mathbb{E}[T_2(N)]=\Theta(N\log N)}.
\]

Algorithm (3) performs $N$ initial assignments and exactly $N-1$
swaps, giving
\[
\boxed{T_3(N)=\Theta(N)}.
\]

\paragraph{(c) C++ implementation.}
The following implementations use the standard Mersenne Twister random
number generator.

\begin{lstlisting}
#include <algorithm>
#include <numeric>
#include <random>
#include <vector>

std::vector<int> permutation1(int n, std::mt19937& rng) {
    std::vector<int> a(n);
    std::uniform_int_distribution<int> pick(1, n);
    for (int i = 0; i < n; ++i) {
        while (true) {
            int value = pick(rng);
            bool duplicate = false;
            for (int j = 0; j < i; ++j)
                if (a[j] == value) {
                    duplicate = true;
                    break;
                }
            if (!duplicate) {
                a[i] = value;
                break;
            }
        }
    }
    return a;
}

std::vector<int> permutation2(int n, std::mt19937& rng) {
    std::vector<int> a(n);
    std::vector<bool> used(n + 1, false);
    std::uniform_int_distribution<int> pick(1, n);
    for (int i = 0; i < n; ++i) {
        int value;
        do {
            value = pick(rng);
        } while (used[value]);
        a[i] = value;
        used[value] = true;
    }
    return a;
}

std::vector<int> permutation3(int n, std::mt19937& rng) {
    std::vector<int> a(n);
    std::iota(a.begin(), a.end(), 1);
    for (int i = 1; i < n; ++i) {
        std::uniform_int_distribution<int> pick(0, i);
        std::swap(a[i], a[pick(rng)]);
    }
    return a;
}
\end{lstlisting}

Each function was executed ten times for each input size. The following
table reports the average elapsed time in milliseconds from a
representative optimized C++17 run.

\begin{center}
\begin{tabular}{c|r|r}
\textbf{Algorithm} & $N$ & \textbf{Average time (ms)} \\
\hline
1 & 250   & 0.069 \\
1 & 500   & 0.307 \\
1 & 1{,}000 & 1.198 \\
1 & 2{,}000 & 4.380 \\
\hline
2 & 25{,}000  & 1.501 \\
2 & 50{,}000  & 2.845 \\
2 & 100{,}000 & 6.255 \\
2 & 200{,}000 & 13.439 \\
2 & 400{,}000 & 28.003 \\
2 & 800{,}000 & 67.784 \\
\hline
3 & 100{,}000   & 0.494 \\
3 & 200{,}000   & 1.035 \\
3 & 400{,}000   & 2.078 \\
3 & 800{,}000   & 4.441 \\
3 & 1{,}600{,}000 & 10.133 \\
3 & 3{,}200{,}000 & 24.167 \\
3 & 6{,}400{,}000 & 76.557
\end{tabular}
\end{center}

\paragraph{(d) Comparison.}
Algorithm (1) becomes expensive at only a few thousand elements, and
its time increases by roughly a factor of four to five when $N$ doubles,
consistent with $\Theta(N^2\log N)$. Algorithm (2) grows a little more
than linearly when $N$ doubles, as predicted by $\Theta(N\log N)$.
Algorithm (3) shows approximately linear scaling until memory-allocation
and cache effects become significant at the largest inputs. The
measurements therefore agree with the expected-time analysis.

\paragraph{(e) Worst-case running times.}
Algorithms (1) and (2) use rejection sampling. An unlucky random-number
sequence can repeatedly generate an already-used value indefinitely, so
neither algorithm has a finite worst-case running-time bound. Their
worst-case running time is unbounded. Algorithm (3) always performs the
same number of assignments and swaps, so its worst-case time is
$\Theta(N)$.

\subsection*{\color{sectionnavy}Exercise 2.9}
\addcontentsline{toc}{subsection}{Exercise 2.9: Extrapolating algorithm runtimes}

\noindent
Complete the table in Figure 2.2 with estimates for the running times
that were too long to simulate. Interpolate the running times for these
algorithms and estimate the time required to compute the maximum
subsequence sum of one million numbers. What assumptions have you made?

\subsubsection*{\color{sectionnavy}Solution}

For a tenfold increase in input size, the expected multiplication
factors are
\[
10^3=1000,\qquad 10^2=100,\qquad
10\left(1+\frac{\log 10}{\log N}\right),\qquad 10
\]
for Algorithms 1 through 4, respectively. Applying these factors to the
most recent available measurements gives the completed table below.
Times are in seconds; the entries marked with an asterisk are estimates.

\begin{center}
\small
\begin{tabular}{r|r|r|r|r}
\textbf{Input size} & \textbf{Alg. 1} & \textbf{Alg. 2} &
\textbf{Alg. 3} & \textbf{Alg. 4} \\
 & $O(N^3)$ & $O(N^2)$ & $O(N\log N)$ & $O(N)$ \\
\hline
100       & 0.000159 & 0.000006 & 0.000005 & 0.000002 \\
1{,}000   & 0.095857 & 0.000371 & 0.000060 & 0.000022 \\
10{,}000  & 86.67    & 0.033322 & 0.000619 & 0.000222 \\
100{,}000 & 86{,}670$^{*}$ & 3.33 & 0.006700 & 0.002205 \\
1{,}000{,}000 & 86{,}670{,}000$^{*}$ & 333$^{*}$ &
0.074870 & 0.022711
\end{tabular}
\end{center}

For one million numbers, the estimated running times are therefore
\[
\begin{array}{c|c|c}
\text{Algorithm} & \text{Seconds} & \text{More readable form} \\
\hline
1 & 86{,}670{,}000 & \text{about }1{,}003\text{ days, or }2.75\text{ years} \\
2 & 333 & \text{about }5.55\text{ minutes} \\
3 & 0.074870 & \text{about }74.9\text{ ms} \\
4 & 0.022711 & \text{about }22.7\text{ ms}
\end{array}
\]

As a check on the interpolation, Algorithm 3 predicted from the
$N=100{,}000$ row would take
\[
0.006700\left(10\frac{\log(1{,}000{,}000)}
                         {\log(100{,}000)}\right)
=0.006700(12)
=0.0804\text{ seconds},
\]
which is close to the measured $0.074870$ seconds. Likewise, the linear
prediction for Algorithm 4 is $10(0.002205)=0.02205$ seconds, close to
the measured $0.022711$ seconds.

The estimates assume that the dominant asymptotic term already controls
the measured running time, its constant factor remains stable, and the
same hardware, compiler, optimization settings, and input conditions are
used. They also ignore secondary terms and effects such as caching,
memory bandwidth, operating-system scheduling, and timer error.

\subsection*{\color{sectionnavy}Exercise 2.10}
\addcontentsline{toc}{subsection}{Exercise 2.10: Hand arithmetic complexity}

\noindent
Determine, for the typical algorithms used to perform calculations by
hand, the running time required to do the following:

\begin{enumerate}
  \renewcommand{\labelenumi}{\alph{enumi}.}
  \item Add two $N$-digit integers.
  \item Multiply two $N$-digit integers.
  \item Divide two $N$-digit integers.
\end{enumerate}

\subsubsection*{\color{sectionnavy}Solution}

Assume that a single-digit arithmetic operation, comparison, or carry
operation takes constant time.

\begin{enumerate}
  \renewcommand{\labelenumi}{\alph{enumi}.}

  \item In ordinary addition, each pair of corresponding digits is
  processed once, together with a possible carry. Thus the running time
  is
  \[
  \boxed{\Theta(N)}.
  \]

  \item In grade-school multiplication, each of the $N$ digits of one
  number is multiplied by each of the $N$ digits of the other. Adding
  the shifted partial products also takes at most quadratic time.
  Therefore, the running time is
  \[
  \boxed{\Theta(N^2)}.
  \]

  \item Long division determines at most $N$ quotient digits. At each
  stage, comparing with and subtracting a suitably shifted divisor can
  require $\Theta(N)$ digit operations. Hence the running time is
  \[
  \boxed{\Theta(N^2)}.
  \]
\end{enumerate}

These bounds describe the usual hand algorithms; faster multiplication
and division algorithms exist for sufficiently large inputs.

\subsection*{\color{sectionnavy}Exercise 2.11}
\addcontentsline{toc}{subsection}{Exercise 2.11: Scaling runtime with input size}

\noindent
An algorithm takes $0.5$ ms for input size $100$. How long will it take
for input size $500$ if the running time is the following (assume
low-order terms are negligible)?

\begin{enumerate}
  \renewcommand{\labelenumi}{\alph{enumi}.}
  \item Linear
  \item $O(N\log N)$
  \item Quadratic
  \item Cubic
\end{enumerate}

\subsubsection*{\color{sectionnavy}Solution}

The input size increases by a factor of
\[
\frac{500}{100}=5.
\]
Scale the original $0.5$ ms running time according to each growth rate.

\begin{enumerate}
  \renewcommand{\labelenumi}{\alph{enumi}.}

  \item For linear time, the multiplication factor is $5$:
  \[
  0.5(5)=\boxed{2.5\text{ ms}}.
  \]

  \item For $N\log N$, the multiplication factor is
  \[
  \frac{500\log 500}{100\log 100}
  =5\frac{\log 500}{\log 100}
  \approx 6.747.
  \]
  Therefore, the estimated time is
  \[
  0.5(6.747)\approx\boxed{3.37\text{ ms}}.
  \]

  \item For quadratic time, the multiplication factor is $5^2=25$:
  \[
  0.5(25)=\boxed{12.5\text{ ms}}.
  \]

  \item For cubic time, the multiplication factor is $5^3=125$:
  \[
  0.5(125)=\boxed{62.5\text{ ms}}.
  \]
\end{enumerate}

\subsection*{\color{sectionnavy}Exercise 2.12}
\addcontentsline{toc}{subsection}{Exercise 2.12: Problem size under a time limit}

\noindent
An algorithm takes $0.5$ ms for input size $100$. How large a problem
can be solved in $1$ min if the running time is the following (assume
low-order terms are negligible)?

\begin{enumerate}
  \renewcommand{\labelenumi}{\alph{enumi}.}
  \item Linear
  \item $O(N\log N)$
  \item Quadratic
  \item Cubic
\end{enumerate}

\subsubsection*{\color{sectionnavy}Solution}

One minute is $60{,}000$ ms, so the available time is larger by the
factor
\[
R=\frac{60{,}000}{0.5}=120{,}000.
\]
If $T(N)=c f(N)$, we therefore solve
\[
\frac{f(N)}{f(100)}=120{,}000.
\]

\begin{enumerate}
  \renewcommand{\labelenumi}{\alph{enumi}.}

  \item For linear time,
  \[
  \frac{N}{100}=120{,}000,
  \qquad
  \boxed{N=12{,}000{,}000}.
  \]

  \item For $N\log N$ time, $N$ satisfies
  \[
  \frac{N\log N}{100\log 100}=120{,}000.
  \]
  Solving this equation numerically gives
  \[
  N\approx3{,}656{,}807,
  \]
  so the largest integral input size is approximately
  \[
  \boxed{N\approx3.66\times10^6}.
  \]

  \item For quadratic time,
  \[
  \left(\frac{N}{100}\right)^2=120{,}000,
  \qquad
  N=100\sqrt{120{,}000}\approx34{,}641.
  \]
  Thus,
  \[
  \boxed{N\approx34{,}641}.
  \]

  \item For cubic time,
  \[
  \left(\frac{N}{100}\right)^3=120{,}000,
  \qquad
  N=100\sqrt[3]{120{,}000}\approx4{,}932.
  \]
  Thus,
  \[
  \boxed{N\approx4{,}932}.
  \]
\end{enumerate}

\subsection*{\color{sectionnavy}Exercise 2.13}
\addcontentsline{toc}{subsection}{Exercise 2.13: Polynomial evaluation by powers}

\noindent
How much time is required to compute
$f(x)=\sum_{i=0}^{N}a_i x^i$?

\begin{enumerate}
  \renewcommand{\labelenumi}{\alph{enumi}.}
  \item Using a simple routine to perform exponentiation?
  \item Using the routine in Section 2.4.4?
\end{enumerate}

\subsubsection*{\color{sectionnavy}Solution}

Assume that each power $x^i$ is computed independently and that each
arithmetic operation takes constant time.

\begin{enumerate}
  \renewcommand{\labelenumi}{\alph{enumi}.}

  \item A simple exponentiation routine computes $x^i$ using
  $i-1=\Theta(i)$ multiplications. Thus the total work is
  \[
  \sum_{i=0}^{N}\Theta(i)=\Theta(N^2).
  \]
  The additional $N+1$ coefficient multiplications and additions do not
  change the bound, so the running time is
  \[
  \boxed{\Theta(N^2)}.
  \]

  \item The fast exponentiation routine from Section 2.4.4 computes
  $x^i$ in $\Theta(\log i)$ time. Computing all powers independently
  therefore requires
  \[
  \sum_{i=1}^{N}\Theta(\log i)
  =\Theta(\log(N!))
  =\Theta(N\log N).
  \]
  Hence the total running time is
  \[
  \boxed{\Theta(N\log N)}.
  \]
\end{enumerate}

\subsection*{\color{sectionnavy}Exercise 2.14}
\addcontentsline{toc}{subsection}{Exercise 2.14: Horner's rule}

\noindent
Consider the following algorithm (known as Horner's rule) to evaluate
$f(x)=\sum_{i=0}^{N}a_i x^i$:

\begin{lstlisting}
poly = 0;
for (i = n; i >= 0; --i)
    poly = x * poly + a[i];
\end{lstlisting}

\begin{enumerate}
  \renewcommand{\labelenumi}{\alph{enumi}.}
  \item Show how the steps are performed by this algorithm for $x=3$,
  $f(x)=4x^4+8x^3+x+2$.
  \item Explain why this algorithm works.
  \item What is the running time of this algorithm?
\end{enumerate}

\subsubsection*{\color{sectionnavy}Solution}

The coefficients in ascending order are
\[
a_0=2,\quad a_1=1,\quad a_2=0,\quad a_3=8,\quad a_4=4.
\]

\begin{enumerate}
  \renewcommand{\labelenumi}{\alph{enumi}.}

  \item Starting with \texttt{poly = 0}, the iterations give
  \[
  \begin{array}{c|c}
  i & \texttt{poly}=3(\texttt{poly})+a_i \\
  \hline
  4 & 3(0)+4=4 \\
  3 & 3(4)+8=20 \\
  2 & 3(20)+0=60 \\
  1 & 3(60)+1=181 \\
  0 & 3(181)+2=545
  \end{array}
  \]
  Therefore,
  \[
  \boxed{f(3)=545}.
  \]

  \item The polynomial can be written in nested form as
  \[
  f(x)=a_0+x\bigl(a_1+x\bigl(a_2+\cdots
       +x(a_{N-1}+xa_N)\cdots\bigr)\bigr).
  \]
  After processing coefficient $a_i$, the variable \texttt{poly}
  contains
  \[
  a_i+a_{i+1}x+a_{i+2}x^2+\cdots+a_Nx^{N-i}.
  \]
  The assignment \texttt{poly = x * poly + a[i]} constructs this
  expression one coefficient at a time, so after the final iteration it
  equals the original polynomial $f(x)$.

  \item The loop executes exactly $N+1$ times, with one multiplication
  and one addition per iteration. Hence its running time is
  \[
  \boxed{\Theta(N)}.
  \]
\end{enumerate}

\subsection*{\color{sectionnavy}Exercise 2.15}
\addcontentsline{toc}{subsection}{Exercise 2.15: Finding a fixed point in an array}

\noindent
Give an efficient algorithm to determine if there exists an integer
$i$ such that $A_i=i$ in an array of integers
$A_1<A_2<A_3<\cdots<A_N$. What is the running time of your algorithm?

\subsubsection*{\color{sectionnavy}Solution}

Use binary search on the index. With the array indexed from $1$ through
$N$, the algorithm is:

\begin{lstlisting}
bool hasFixedPoint(const vector<int>& a) {
    int low = 1;
    int high = static_cast<int>(a.size()) - 1;

    while (low <= high) {
        int mid = low + (high - low) / 2;
        if (a[mid] == mid)
            return true;
        if (a[mid] < mid)
            low = mid + 1;
        else
            high = mid - 1;
    }
    return false;
}
\end{lstlisting}

To justify the discarded halves, suppose first that $A_{\mathit{mid}}<
\mathit{mid}$. Because the entries are strictly increasing integers,
for every $j<\mathit{mid}$,
\[
A_j\leq A_{\mathit{mid}}-(\mathit{mid}-j)<j.
\]
Thus no index to the left can satisfy $A_j=j$. Similarly, if
$A_{\mathit{mid}}>\mathit{mid}$, then for every $j>\mathit{mid}$,
\[
A_j\geq A_{\mathit{mid}}+(j-\mathit{mid})>j,
\]
so no index to the right can be a solution.

Each comparison eliminates at least half of the remaining indices.
Therefore, the running time is
\[
\boxed{\Theta(\log N)},
\]
with $\Theta(1)$ auxiliary space for the iterative implementation.

\subsection*{\color{sectionnavy}Exercise 2.16}
\addcontentsline{toc}{subsection}{Exercise 2.16: The binary GCD algorithm}

\noindent
Write an alternative gcd algorithm based on the following observations
(arrange so that $a>b$):

\begin{enumerate}
  \renewcommand{\labelenumi}{\alph{enumi}.}
  \item $\gcd(a,b)=2\gcd(a/2,b/2)$ if $a$ and $b$ are both even.
  \item $\gcd(a,b)=\gcd(a/2,b)$ if $a$ is even and $b$ is odd.
  \item $\gcd(a,b)=\gcd(a,b/2)$ if $a$ is odd and $b$ is even.
  \item $\gcd(a,b)=\gcd((a+b)/2,(a-b)/2)$ if $a$ and $b$ are both
  odd.
\end{enumerate}

\subsubsection*{\color{sectionnavy}Solution}

The following recursive implementation is a form of the binary gcd
algorithm. The expression used for $(a+b)/2$ avoids overflow from
forming $a+b$ directly.

\begin{lstlisting}
#include <cstdint>
#include <utility>

std::uint64_t binaryGcd(std::uint64_t a, std::uint64_t b) {
    if (a == 0) return b;
    if (b == 0) return a;
    if (a < b) std::swap(a, b);
    if (a == b) return a;

    bool aEven = (a % 2 == 0);
    bool bEven = (b % 2 == 0);

    if (aEven && bEven)
        return 2 * binaryGcd(a / 2, b / 2);
    if (aEven)
        return binaryGcd(a / 2, b);
    if (bEven)
        return binaryGcd(a, b / 2);

    // Both are odd and a > b.
    return binaryGcd(b + (a - b) / 2, (a - b) / 2);
}
\end{lstlisting}

The first identity removes a common factor of $2$. In the second and
third cases, an odd number cannot share a factor of $2$ with the even
number, so dividing the even argument by $2$ does not change the gcd.

For the final case, define
\[
u=\frac{a+b}{2},\qquad v=\frac{a-b}{2}.
\]
Then $a=u+v$ and $b=u-v$. Consequently, every common divisor of $u$
and $v$ divides both $a$ and $b$. Conversely, since $a$ and $b$ are
odd, their gcd is odd and divides both $(a+b)/2$ and $(a-b)/2$.
Therefore $\gcd(a,b)=\gcd(u,v)$.

Each call either halves an even argument or replaces two odd arguments
by values formed from their half-sum and half-difference. After at most
a constant number of calls, the magnitude of the arguments decreases
by a constant factor. Under the usual unit-cost arithmetic model, the
worst-case running time is therefore
\[
\boxed{\Theta(\log \max(a,b))}.
\]
The recursion also uses $\Theta(\log \max(a,b))$ stack space; an
iterative version can use constant auxiliary space.

\subsection*{\color{sectionnavy}Exercise 2.17}
\addcontentsline{toc}{subsection}{Exercise 2.17: Subsequence sums and products}

\noindent
Give efficient algorithms (along with running-time analyses) to:

\begin{enumerate}
  \renewcommand{\labelenumi}{\alph{enumi}.}
  \item Find the minimum subsequence sum.
  \item Find the minimum positive subsequence sum.
  \item Find the maximum subsequence product.
\end{enumerate}

\subsubsection*{\color{sectionnavy}Solution}

Here, a subsequence means a contiguous subsequence.

\begin{enumerate}
  \renewcommand{\labelenumi}{\alph{enumi}.}

  \item This is the minimum-sum version of Kadane's algorithm. If the
  empty subsequence is permitted, initialize both values to zero.

\begin{lstlisting}
long long minimumSubsequenceSum(const vector<int>& a) {
    long long endingHere = 0;
    long long minimumSum = 0;

    for (int x : a) {
        endingHere = std::min(0LL, endingHere + x);
        minimumSum = std::min(minimumSum, endingHere);
    }
    return minimumSum;
}
\end{lstlisting}

  Each element is processed once, so the running time is
  $\boxed{\Theta(N)}$ and the auxiliary space is $\Theta(1)$.

  \item Let the prefix sums be
  \[
  P_0=0,\qquad P_j=\sum_{k=1}^{j}A_k.
  \]
  A subsequence from $i+1$ through $j$ has sum $P_j-P_i$. For each
  $P_j$, the smallest positive difference is obtained by choosing the
  largest earlier prefix sum that is strictly less than $P_j$.
  Maintain all earlier prefix sums in an ordered set and perform a
  predecessor query for every $P_j$.

\begin{lstlisting}
long long minimumPositiveSubsequenceSum(const vector<int>& a) {
    std::set<long long> earlier{0};
    long long prefix = 0;
    long long answer = std::numeric_limits<long long>::max();

    for (int x : a) {
        prefix += x;
        auto it = earlier.lower_bound(prefix);
        if (it != earlier.begin()) {
            --it;
            answer = std::min(answer, prefix - *it);
        }
        earlier.insert(prefix);
    }
    return answer; // No answer if this remains LLONG_MAX.
}
\end{lstlisting}

  Each ordered-set operation takes $O(\log N)$ time, so the total is
  $\boxed{O(N\log N)}$ time and $O(N)$ space.

  \item A negative value can turn the smallest product ending at the
  previous position into the largest. Therefore, keep both the maximum
  and minimum products ending at the current position.

\begin{lstlisting}
long long maximumSubsequenceProduct(const vector<int>& a) {
    long long maximumEnding = a[0];
    long long minimumEnding = a[0];
    long long answer = a[0];

    for (std::size_t i = 1; i < a.size(); ++i) {
        long long x = a[i];
        if (x < 0)
            std::swap(maximumEnding, minimumEnding);

        maximumEnding = std::max(x, x * maximumEnding);
        minimumEnding = std::min(x, x * minimumEnding);
        answer = std::max(answer, maximumEnding);
    }
    return answer;
}
\end{lstlisting}

  This algorithm takes $\boxed{\Theta(N)}$ time and $\Theta(1)$
  auxiliary space. The numeric type must be large enough to avoid
  overflow.
\end{enumerate}

\subsection*{\color{sectionnavy}Exercise 2.18}
\addcontentsline{toc}{subsection}{Exercise 2.18: Root finding by bisection}

\noindent
An important problem in numerical analysis is to find a solution to
$f(x)=0$ for an arbitrary continuous function $f$. If $f(\mathit{low})$
and $f(\mathit{high})$ have opposite signs, write a function that uses
binary search to solve for a zero. What must be done to ensure
termination?

\subsubsection*{\color{sectionnavy}Solution}

The bisection method repeatedly keeps the half-interval whose endpoint
values have opposite signs.

\begin{lstlisting}
double findZero(const std::function<double(double)>& f,
                double low, double high,
                double xTolerance = 1e-12,
                double fTolerance = 1e-12) {
    double fLow = f(low);
    double fHigh = f(high);

    if (std::abs(fLow) <= fTolerance) return low;
    if (std::abs(fHigh) <= fTolerance) return high;
    if (std::signbit(fLow) == std::signbit(fHigh))
        throw std::invalid_argument("root is not bracketed");

    for (int iteration = 0; iteration < 200; ++iteration) {
        double mid = low + (high - low) / 2.0;
        double fMid = f(mid);

        if (std::abs(fMid) <= fTolerance ||
            high - low <= xTolerance ||
            mid == low || mid == high)
            return mid;

        if (std::signbit(fLow) != std::signbit(fMid)) {
            high = mid;
        } else {
            low = mid;
            fLow = fMid;
        }
    }
    return low + (high - low) / 2.0;
}
\end{lstlisting}

Testing only whether $f(\mathit{mid})=0$ does not guarantee termination,
because a floating-point computation may never produce the exact root.
Termination must instead be based on a tolerance for the interval or
function value, together with a maximum iteration count. After $k$
iterations the interval width is $(\mathit{high}-\mathit{low})/2^k$, so
an $x$-accuracy of $\varepsilon$ requires
\[
\boxed{O\!\left(\log\frac{\mathit{high}-\mathit{low}}{\varepsilon}
       \right)}
\]
function evaluations.

\subsection*{\color{sectionnavy}Exercise 2.19}
\addcontentsline{toc}{subsection}{Exercise 2.19: Returning subsequence endpoints}

\noindent
Modify the maximum contiguous subsequence sum algorithms so that they
return, in a single object, both the maximum sum and the indices of the
actual subsequence.

\subsubsection*{\color{sectionnavy}Solution}

Store the sum and its two endpoints in a result object. The linear-time
algorithm becomes:

\begin{lstlisting}
struct SubsequenceResult {
    long long sum;
    int first;
    int last;
};

SubsequenceResult maximumSubsequence(const vector<int>& a) {
    SubsequenceResult best{0, -1, -1};
    long long currentSum = 0;
    int currentFirst = 0;

    for (int i = 0; i < static_cast<int>(a.size()); ++i) {
        currentSum += a[i];

        if (currentSum > best.sum)
            best = {currentSum, currentFirst, i};

        if (currentSum < 0) {
            currentSum = 0;
            currentFirst = i + 1;
        }
    }
    return best;
}
\end{lstlisting}

The indices $(-1,-1)$ represent the empty subsequence when all entries
are negative. The running time remains $\boxed{\Theta(N)}$, since only
constant additional work is performed per element.

The same idea applies to the other algorithms in the text. In the
cubic and quadratic versions, record the current $i$ and $j$ whenever a
new maximum is found. In the divide-and-conquer version, each recursive
call returns a \texttt{SubsequenceResult}; the left, right, and crossing
results are constructed with their endpoints, and the one with the
largest sum is returned. These additions do not change their respective
asymptotic running times.

\subsection*{\color{sectionnavy}Exercise 2.20}
\addcontentsline{toc}{subsection}{Exercise 2.20: Trial division and input bits}

\noindent
Consider a program that determines whether a positive integer $N$ is
prime.

\begin{enumerate}
  \renewcommand{\labelenumi}{\alph{enumi}.}
  \item Write the program.
  \item Give its worst-case running time in terms of $N$.
  \item Let $B$ be the number of bits in the binary representation of
  $N$. What is $B$?
  \item Give the worst-case running time in terms of $B$.
  \item Compare testing a 20-bit number and a 40-bit number.
  \item Is it more reasonable to state the running time in terms of
  $N$ or $B$? Why?
\end{enumerate}

\subsubsection*{\color{sectionnavy}Solution}

\begin{enumerate}
  \renewcommand{\labelenumi}{\alph{enumi}.}

  \item It is sufficient to test possible divisors through $\sqrt N$.
  The comparison $d\leq N/d$ avoids overflow from computing $d^2$.

\begin{lstlisting}
bool isPrime(std::uint64_t n) {
    if (n < 2) return false;
    if (n % 2 == 0) return n == 2;

    for (std::uint64_t d = 3; d <= n / d; d += 2)
        if (n % d == 0)
            return false;
    return true;
}
\end{lstlisting}

  If $N$ is composite, it has a factor no larger than $\sqrt N$; if no
  such factor is found, $N$ is prime.

  \item In the worst case, such as when $N$ is prime, the program tests
  $\Theta(\sqrt N)$ possible divisors. Under the unit-cost arithmetic
  assumption, its running time is
  \[
  \boxed{\Theta(\sqrt N)}.
  \]

  \item For $N\geq1$, the binary input length is
  \[
  \boxed{B=\lfloor\log_2N\rfloor+1}.
  \]

  \item Since $N$ can be approximately $2^B$,
  \[
  \sqrt N=\Theta(2^{B/2}).
  \]
  Thus the unit-cost running time in terms of the input length is
  \[
  \boxed{\Theta(2^{B/2})},
  \]
  which is exponential in $B$. Accounting for the bit cost of division
  would add a polynomial factor in $B$.

  \item A 20-bit number has square root on the order of $2^{10}$,
  whereas a 40-bit number has square root on the order of $2^{20}$.
  Therefore the 40-bit test can require approximately
  \[
  \frac{2^{20}}{2^{10}}=\boxed{1024}
  \]
  times as many trial divisions.

  \item It is more informative to express the running time in terms of
  $B$, because $B$ is the actual size of the input representation. The
  bound $\Theta(\sqrt N)$ can appear moderate, while the equivalent
  $\Theta(2^{B/2})$ correctly reveals that trial division is exponential
  in the number of input bits.
\end{enumerate}

\subsection*{\color{sectionnavy}Exercise 2.21}
\addcontentsline{toc}{subsection}{Exercise 2.21: Sieve of Eratosthenes complexity}

\noindent
The Sieve of Eratosthenes computes all primes less than $N$ by crossing
out multiples of each uncrossed integer through $\sqrt N$. What is its
running time?

\subsubsection*{\color{sectionnavy}Solution}

Initializing and scanning the Boolean table takes $\Theta(N)$ time. For
each prime $p\leq\sqrt N$, crossing out its multiples performs at most
$N/p$ operations. Hence the total marking work is
\[
\sum_{\substack{p\leq\sqrt N\\p\text{ prime}}}\frac{N}{p}
=N\sum_{p\leq\sqrt N}\frac1p.
\]
The sum of the reciprocals of the primes up to $x$ is
$\Theta(\log\log x)$, so
\[
N\sum_{p\leq\sqrt N}\frac1p
=\Theta(N\log\log N).
\]
The linear initialization and scan are dominated by this term.
Therefore, assuming constant-time table operations, the running time is
\[
\boxed{\Theta(N\log\log N)},
\]
and the table requires $\boxed{\Theta(N)}$ space.

\subsection*{\color{sectionnavy}Exercise 2.22}
\addcontentsline{toc}{subsection}{Exercise 2.22: Computing a power efficiently}

\noindent
Show that $X^{62}$ can be computed with only eight multiplications.

\subsubsection*{\color{sectionnavy}Solution}

Use the following addition chain for the exponents:
\[
1,\ 2,\ 3,\ 5,\ 10,\ 20,\ 30,\ 31,\ 62.
\]
It gives the computation
\[
\begin{aligned}
X^2   &=X\cdot X,                 &
X^3   &=X^2\cdot X,\\
X^5   &=X^3\cdot X^2,             &
X^{10}&=X^5\cdot X^5,\\
X^{20}&=X^{10}\cdot X^{10},       &
X^{30}&=X^{20}\cdot X^{10},\\
X^{31}&=X^{30}\cdot X,            &
X^{62}&=X^{31}\cdot X^{31}.
\end{aligned}
\]
Each line contains one multiplication, for a total of exactly
\[
\boxed{8\text{ multiplications}}.
\]

\subsection*{\color{sectionnavy}Exercise 2.23}
\addcontentsline{toc}{subsection}{Exercise 2.23: Iterative fast exponentiation}

\noindent
Write the fast exponentiation routine without recursion.

\subsubsection*{\color{sectionnavy}Solution}

The exponent can be processed from its most significant binary digit to
its least significant digit. The leading digit is always $1$, so the
initial value is $x$.

\begin{lstlisting}
long long power(long long x, unsigned int n) {
    if (n == 0)
        return 1;

    unsigned int mask = 1;
    while (mask <= n / 2)
        mask <<= 1;              // Highest power of two <= n.

    long long result = x;        // Process the leading 1 bit.
    mask >>= 1;

    while (mask != 0) {
        result *= result;
        if ((n & mask) != 0)
            result *= x;
        mask >>= 1;
    }
    return result;
}
\end{lstlisting}

Squaring shifts the already processed binary exponent one position to
the left; multiplying by $x$ appends a $1$ bit. Since the number of bits
in $n$ is $\Theta(\log n)$, the running time is
\[
\boxed{\Theta(\log n)}.
\]

\subsection*{\color{sectionnavy}Exercise 2.24}
\addcontentsline{toc}{subsection}{Exercise 2.24: Counting exponentiation multiplications}

\noindent
Give a precise count of the number of multiplications used by the fast
exponentiation routine.

\subsubsection*{\color{sectionnavy}Solution}

Let
\[
n=(b_kb_{k-1}\cdots b_1b_0)_2,
\qquad b_k=1,
\]
and let $s(n)$ denote the number of $1$ bits in this binary
representation. There are
\[
k=\lfloor\log_2 n\rfloor
\]
bits after the leading bit. Processing each of them requires one
squaring, giving $k$ multiplications. Every remaining $1$ bit also
requires one multiplication by $x$, giving $s(n)-1$ additional
multiplications. Thus, for $n\geq1$, the precise count is
\[
\boxed{M(n)=\lfloor\log_2 n\rfloor+s(n)-1}.
\]
For $n=0$, no multiplication is required. Since
$1\leq s(n)\leq\lfloor\log_2n\rfloor+1$, this count lies between
\[
\lfloor\log_2n\rfloor
\quad\text{and}\quad
2\lfloor\log_2n\rfloor,
\]
confirming the $\Theta(\log n)$ bound. For example, $62=(111110)_2$,
so this routine uses $5+5-1=9$ multiplications; the specially chosen
addition chain in Exercise 2.22 improves this to eight.

\subsection*{\color{sectionnavy}Exercise 2.25}
\addcontentsline{toc}{subsection}{Exercise 2.25: Comparing runtime guarantees}

\noindent
Programs $A$ and $B$ have worst-case running times no greater than
$150N\log_2N$ and $N^2$, respectively. Answer the following questions,
if possible.

\begin{enumerate}
  \renewcommand{\labelenumi}{\alph{enumi}.}
  \item Which program has the better guarantee for large values
  ($N>10{,}000$)?
  \item Which program has the better guarantee for small values
  ($N<100$)?
  \item Which program will run faster on average for $N=1{,}000$?
  \item Is it possible that program $B$ runs faster than program $A$ on
  all possible inputs?
\end{enumerate}

\subsubsection*{\color{sectionnavy}Solution}

Compare the two stated upper bounds by cancelling the positive factor
$N$:
\[
150N\log_2N<N^2
\quad\Longleftrightarrow\quad
150\log_2N<N.
\]

\begin{enumerate}
  \renewcommand{\labelenumi}{\alph{enumi}.}

  \item At $N=10{,}000$,
  \[
  150\log_2(10{,}000)\approx1{,}993<10{,}000,
  \]
  and the linear factor $N$ eventually grows faster than $\log N$.
  Thus $150N\log_2N<N^2$ throughout the stated large-input range, so
  \[
  \boxed{\text{Program }A\text{ has the better worst-case guarantee.}}
  \]

  \item For the relevant integers $2\leq N<100$,
  $150\log_2N>N$, and hence $N^2<150N\log_2N$. Therefore,
  \[
  \boxed{\text{Program }B\text{ has the better worst-case guarantee.}}
  \]

  \item This cannot be determined. The given information describes
  only worst-case upper bounds; it provides neither average-case running
  times nor an input distribution. Thus
  \[
  \boxed{\text{there is insufficient information.}}
  \]

  \item Yes. An upper bound does not state that a program actually takes
  that much time, nor does it provide a lower bound. Program $B$ could
  have a smaller actual running time on every input while still obeying
  its $N^2$ upper bound. Hence
  \[
  \boxed{\text{yes, this is possible.}}
  \]
\end{enumerate}

\subsection*{\color{sectionnavy}Exercise 2.26}
\addcontentsline{toc}{subsection}{Exercise 2.26: Finding a majority element}

\noindent
A majority element in an array $A$ of size $N$ is an element that
appears more than $N/2$ times. To find a candidate, compare adjacent
pairs, place one copy of each equal pair in a second array $B$, discard
unequal pairs, and recursively find a candidate in $B$. Then scan $A$
to verify whether the candidate is actually a majority.

\begin{enumerate}
  \renewcommand{\labelenumi}{\alph{enumi}.}
  \item How does the recursion terminate?
  \item How is the case where $N$ is odd handled?
  \item What is the running time of the algorithm?
  \item How can the extra array $B$ be avoided?
  \item Write a program to compute the majority element.
\end{enumerate}

\subsubsection*{\color{sectionnavy}Solution}

Pair cancellation is valid because discarding two unequal elements
removes at most one occurrence of any particular value. Thus, if an
element is a majority before the cancellation, it remains the only
possible majority candidate afterward.

\begin{enumerate}
  \renewcommand{\labelenumi}{\alph{enumi}.}

  \item The recursion stops when the current array has zero or one
  element. An empty array has no candidate, while the sole element of a
  one-element array is its candidate. The recursion also stops if a
  pairing step produces an empty array.

  \item Let $N=2k+1$ and temporarily leave the final element $z$
  unpaired. Form $B$ from the first $2k$ elements and recursively obtain
  its candidate, if one exists. There are then at most two candidates
  that must be checked in the original array: the candidate from $B$
  and $z$.

  To see why, suppose the true majority is $x$. If $z\neq x$, then all
  of the more than $N/2$ copies of $x$ occur in the first $2k$
  positions, where $x$ is still a majority and must survive as the
  recursive candidate. If $z=x$, the remaining copies can be exactly
  half of the first $2k$ entries, so the recursive reduction might
  discard them all; in that case $z$ supplies the missing candidate.

  \item Forming the paired array takes $\Theta(N)$ time and leaves at
  most $N/2$ elements. Candidate verification is another linear scan,
  so the recurrence is
  \[
  T(N)=T(N/2)+\Theta(N).
  \]
  Expanding it gives
  \[
  T(N)=\Theta\left(N+\frac N2+\frac N4+\cdots\right)
      =\boxed{\Theta(N)}.
  \]
  A linear lower bound is unavoidable because an element near the end
  can change whether a candidate has a strict majority.

  \item The equal-pair representatives can be written into the
  beginning of $A$ itself. Maintain a write index, scan the current
  prefix two elements at a time, and copy one member of every equal pair
  to \texttt{A[write++]}. The recursive call then operates only on this
  shortened prefix. This removes the auxiliary array, although the
  recursive version still uses $O(\log N)$ stack space.

  An even simpler constant-space implementation uses the equivalent
  Boyer--Moore cancellation process: keep one candidate and a balance,
  cancelling each occurrence of a different value against one
  occurrence of the candidate.

  \item The candidate produced by cancellation must be verified because
  an array need not contain any majority element.

\begin{lstlisting}
#include <optional>
#include <vector>

std::optional<int> majorityElement(const std::vector<int>& a) {
    if (a.empty())
        return std::nullopt;

    int candidate = 0;
    int balance = 0;

    for (int value : a) {
        if (balance == 0) {
            candidate = value;
            balance = 1;
        } else if (value == candidate) {
            ++balance;
        } else {
            --balance;
        }
    }

    std::size_t occurrences = 0;
    for (int value : a)
        if (value == candidate)
            ++occurrences;

    if (occurrences > a.size() / 2)
        return candidate;
    return std::nullopt;
}
\end{lstlisting}

  The first scan finds the only possible candidate and the second scan
  verifies it. The program therefore uses $\boxed{\Theta(N)}$ time and
  $\boxed{\Theta(1)}$ auxiliary space.
\end{enumerate}

\subsection*{\color{sectionnavy}Exercise 2.27}
\addcontentsline{toc}{subsection}{Exercise 2.27: Searching a sorted matrix}

\noindent
The input is an $N$ by $N$ matrix already in memory. Every row is
increasing from left to right, and every column is increasing from top
to bottom. Give an $O(N)$ worst-case algorithm that determines whether
a number $X$ occurs in the matrix.

\subsubsection*{\color{sectionnavy}Solution}

Begin at the upper-right corner. If the current element is larger than
$X$, move left; if it is smaller than $X$, move down.

\begin{lstlisting}
bool contains(const vector<vector<int>>& matrix, int x) {
    int n = static_cast<int>(matrix.size());
    int row = 0;
    int column = n - 1;

    while (row < n && column >= 0) {
        if (matrix[row][column] == x)
            return true;
        if (matrix[row][column] > x)
            --column;
        else
            ++row;
    }
    return false;
}
\end{lstlisting}

If the current value is greater than $X$, every element below it in the
same column is still larger, so that column can be discarded. If the
current value is less than $X$, every element to its left in the same
row is still smaller, so that row can be discarded. Thus every move
eliminates one complete row or column without discarding a possible
match.

There are at most $N$ downward moves and $N$ leftward moves, so the
algorithm performs at most $2N-1$ comparisons. Its worst-case running
time is therefore
\[
\boxed{\Theta(N)},
\]
with $\Theta(1)$ auxiliary space.

\subsection*{\color{sectionnavy}Exercise 2.28}
\addcontentsline{toc}{subsection}{Exercise 2.28: Optimizing ordered array pairs}

\noindent
Design efficient algorithms that take an array of positive numbers
$a$ and determine, subject to $j\geq i$, the maximum value of:

\begin{enumerate}
  \renewcommand{\labelenumi}{\alph{enumi}.}
  \item $a[j]+a[i]$
  \item $a[j]-a[i]$
  \item $a[j]a[i]$
  \item $a[j]/a[i]$
\end{enumerate}

\subsubsection*{\color{sectionnavy}Solution}

For each position $j$, maintain the minimum and maximum values in the
prefix $a[0],\ldots,a[j]$. This ensures that the chosen index $i$ always
satisfies $i\leq j$.

\begin{lstlisting}
struct PairExtrema {
    double maximumSum;
    double maximumDifference;
    double maximumProduct;
    double maximumRatio;
};

PairExtrema findExtrema(const vector<double>& a) {
    double prefixMinimum = a[0];
    double prefixMaximum = a[0];

    PairExtrema answer{
        2 * a[0], 0.0, a[0] * a[0], 1.0
    };

    for (double value : a) {
        prefixMinimum = std::min(prefixMinimum, value);
        prefixMaximum = std::max(prefixMaximum, value);

        answer.maximumSum = std::max(
            answer.maximumSum, value + prefixMaximum);
        answer.maximumDifference = std::max(
            answer.maximumDifference, value - prefixMinimum);
        answer.maximumProduct = std::max(
            answer.maximumProduct, value * prefixMaximum);
        answer.maximumRatio = std::max(
            answer.maximumRatio, value / prefixMinimum);
    }
    return answer;
}
\end{lstlisting}

Because $i=j$ is allowed and all values are positive, parts (a) and
(c) simplify to
\[
\boxed{2\max_k a[k]}
\qquad\text{and}\qquad
\boxed{\left(\max_k a[k]\right)^2},
\]
respectively. For each fixed $j$, the best choice in part (b) is the
smallest earlier value, giving
\[
\boxed{\max_j\left(a[j]-\min_{i\leq j}a[i]\right)},
\]
and positivity makes the same prefix minimum the best denominator in
part (d):
\[
\boxed{\max_j\left(\frac{a[j]}{\min_{i\leq j}a[i]}\right)}.
\]
All four answers are obtained in one pass, so the running time is
$\boxed{\Theta(N)}$ and the auxiliary space is $\boxed{\Theta(1)}$.

\subsection*{\color{sectionnavy}Exercise 2.29}
\addcontentsline{toc}{subsection}{Exercise 2.29: Fixed-size integer assumptions}

\noindent
Why is it important to assume that integers in our computer model have
a fixed size?

\subsubsection*{\color{sectionnavy}Solution}

The usual RAM model counts an integer arithmetic operation, comparison,
assignment, or memory access as taking constant time. This is realistic
only when every integer fits in a fixed-size machine word.

If integers could contain arbitrarily many bits while addition,
multiplication, division, and comparison were still charged constant
time, a single ``integer operation'' could process an unbounded amount
of information. For example, an entire long input could be encoded in
one enormous integer, and the model would incorrectly treat operations
on that encoding as taking the same time as operations on a small
integer.

With arbitrary-precision integers, the cost must instead depend on the
operand length. Addition of two $B$-bit integers takes $\Theta(B)$ time
with the ordinary method, while multiplication and division require
still more work. The input size of a numerical value $N$ is also
$\Theta(\log N)$ bits rather than $N$ itself.

Therefore, the fixed-size assumption keeps constant-time primitive
operations and the resulting running-time analyses meaningful. When
values can grow beyond a machine word, their bit lengths must be
included explicitly in the analysis.

\subsection*{\color{sectionnavy}Exercise 2.30}
\addcontentsline{toc}{subsection}{Exercise 2.30: Word-puzzle search complexity}

\noindent
Consider the word-puzzle problem from Chapter 1. Suppose the length of
the longest word is fixed at 10 characters.

\begin{enumerate}
  \renewcommand{\labelenumi}{\alph{enumi}.}
  \item In terms of $R$ and $C$, the numbers of rows and columns, and
  $W$, the number of words, what are the running times of the algorithms
  described in Chapter 1?
  \item Suppose the word list is presorted. Show how binary search gives
  an algorithm with significantly better running time.
\end{enumerate}

\subsubsection*{\color{sectionnavy}Solution}

There are eight possible directions from any starting cell. Since the
maximum word length is the fixed constant 10, checking or constructing
one possible word takes $O(1)$ time.

\begin{enumerate}
  \renewcommand{\labelenumi}{\alph{enumi}.}

  \item The first algorithm considers each of the $W$ dictionary words
  at every one of the $RC$ starting cells and in each of the eight
  directions. A comparison examines at most 10 characters. Therefore,
  \[
  T_1(R,C,W)=O(8\cdot10\cdot RCW)
            =\boxed{O(RCW)}.
  \]

  The second algorithm considers each starting cell, each direction,
  and each possible length from 1 through 10. This produces only
  $O(8\cdot10\cdot RC)=O(RC)$ candidate strings. Testing a candidate by
  a sequential scan of the $W$-word list takes $O(W)$ time, so
  \[
  T_2(R,C,W)=O(RC\cdot W)=\boxed{O(RCW)}.
  \]

  Thus the two straightforward algorithms have the same asymptotic
  worst-case bound when the maximum word length is fixed, although their
  constant factors can differ.

  \item Use the second algorithm, but look up each generated candidate
  string with binary search in the sorted word list. A lookup now takes
  $O(\log W)$ comparisons rather than $O(W)$ comparisons.

\begin{lstlisting}
for each cell (row, column)
    for each of the 8 directions
        word = ""
        for length = 1 to 10
            if the next cell is outside the board
                break
            append the next character to word
            if binarySearch(dictionary, word)
                report word
\end{lstlisting}

  There are $O(RC)$ generated strings because both the number of
  directions and the maximum length are constants. Consequently, the
  running time becomes
  \[
  \boxed{O(RC\log W)}.
  \]
  The word list itself can be sorted once in $O(W\log W)$ time if it is
  not already sorted. If the same dictionary is used for many puzzles,
  this preprocessing cost is paid only once.
\end{enumerate}

\subsection*{\color{sectionnavy}Exercise 2.31}
\addcontentsline{toc}{subsection}{Exercise 2.31: A binary-search termination bug}

\noindent
Suppose that the binary-search routine used \texttt{low = mid} instead
of \texttt{low = mid + 1}. Would the routine still work?

\subsubsection*{\color{sectionnavy}Solution}

No. When the search interval contains two adjacent positions,
\[
\mathit{high}=\mathit{low}+1,
\]
integer division gives
\[
\mathit{mid}
=\left\lfloor\frac{\mathit{low}+\mathit{high}}2\right\rfloor
=\mathit{low}.
\]
If $a[\mathit{mid}]<x$, the assignment
\texttt{low = mid} leaves \texttt{low}, \texttt{mid}, and
\texttt{high} unchanged. The next iteration repeats exactly the same
comparison, so the loop does not terminate.

For example, searching for $4$ in $[1,3]$ first examines the element
$1$. Setting \texttt{low = mid} keeps the interval at $[0,1]$ forever.
The correct assignment is
\[
\boxed{\texttt{low = mid + 1}},
\]
because the comparison has already proved that position
\texttt{mid} cannot contain the target.

\subsection*{\color{sectionnavy}Exercise 2.32}
\addcontentsline{toc}{subsection}{Exercise 2.32: Binary search with one comparison}

\noindent
Implement binary search so that only one two-way comparison is
performed in each iteration.

\subsubsection*{\color{sectionnavy}Solution}

Maintain the half-open interval $[\mathit{low},\mathit{high})$ and find
the first position whose value is not less than $x$. Each iteration asks
only whether \texttt{a[mid] < x}.

\begin{lstlisting}
template <typename Comparable>
int binarySearch(const vector<Comparable>& a,
                 const Comparable& x) {
    int low = 0;
    int high = static_cast<int>(a.size());

    while (low < high) {
        int mid = low + (high - low) / 2;

        if (a[mid] < x)          // The only comparison in the loop.
            low = mid + 1;
        else
            high = mid;
    }

    if (low < static_cast<int>(a.size()) && a[low] == x)
        return low;
    return -1;
}
\end{lstlisting}

On termination, \texttt{low} is the first index at which an occurrence
of $x$ could appear. One final equality test determines whether $x$ is
actually present. The running time remains
\[
\boxed{\Theta(\log N)}.
\]

\subsection*{\color{sectionnavy}Exercise 2.33}
\addcontentsline{toc}{subsection}{Exercise 2.33: Incorrect recursive partitioning}

\noindent
Suppose that the recursive calls in the divide-and-conquer maximum
subsequence sum algorithm are changed to

\begin{lstlisting}
int maxLeftSum  = maxSumRec(a, left, center - 1);
int maxRightSum = maxSumRec(a, center, right);
\end{lstlisting}

Would the routine still work?

\subsubsection*{\color{sectionnavy}Solution}

No. Consider a subproblem containing exactly two elements, so that
\[
\mathit{right}=\mathit{left}+1.
\]
With the original midpoint calculation,
\[
\mathit{center}
=\left\lfloor\frac{\mathit{left}+\mathit{right}}2\right\rfloor
=\mathit{left}.
\]
The proposed right recursive call is then
\[
\operatorname{maxSumRec}(a,\mathit{center},\mathit{right})
=\operatorname{maxSumRec}(a,\mathit{left},\mathit{right}),
\]
which is identical to the current call. The subproblem does not become
smaller, so the routine recurses indefinitely and eventually exhausts
the call stack.

The original partition
\[
[\mathit{left},\mathit{center}]
\quad\text{and}\quad
[\mathit{center}+1,\mathit{right}]
\]
guarantees that both recursive subproblems are strictly smaller. Thus
the proposed replacement
\[
\boxed{\text{does not work with the original midpoint and base case.}}
\]

\subsection*{\color{sectionnavy}Exercise 2.34}
\addcontentsline{toc}{subsection}{Exercise 2.34: Counting nested-loop executions}

\noindent
The innermost code of the cubic maximum subsequence sum algorithm runs
$N(N+1)(N+2)/6$ times, the quadratic version runs $N(N+1)/2$ times,
and the linear version runs $N$ times. What pattern is evident, and what
is its combinatorial explanation?

\subsubsection*{\color{sectionnavy}Solution}

The three counts are consecutive combinations with repetition:
\[
N=\binom{N}{1},
\qquad
\frac{N(N+1)}2=\binom{N+1}{2},
\qquad
\frac{N(N+1)(N+2)}6=\binom{N+2}{3}.
\]
Thus the evident pattern is
\[
\boxed{\binom{N+k-1}{k}}
\]
for $k$ nested indices constrained to occur in nondecreasing order.

In the quadratic algorithm, an execution corresponds to a pair of
indices
\[
0\leq i\leq j<N.
\]
This is the same as choosing a multiset of two indices from $N$
possible positions, of which there are $\binom{N+1}{2}$.

In the cubic algorithm, an execution of the innermost statement is
identified by a triple
\[
0\leq i\leq k\leq j<N,
\]
where $i$ and $j$ are the subsequence endpoints and $k$ is the element
being added. Such triples are multisets of three indices selected from
$N$ positions, so their number is
\[
\binom{N+2}{3}=\frac{N(N+1)(N+2)}6.
\]
The linear algorithm selects only one index, giving
$\binom{N}{1}=N$. This multiset interpretation explains all three
formulas.

\end{document}
